\documentclass[../DoAn.tex]{subfiles}
\begin{document}

Chương này giới thiệu các công nghệ và nền tảng lý thuyết được sử dụng trong hệ thống. Với mỗi công nghệ, chương trình bày nguyên lý hoạt động cơ bản, liệt kê các lựa chọn thay thế, và phân tích lý do lựa chọn dựa trên các yêu cầu đã xác định ở Chương 2.

\section{Phần cứng nhúng}
\label{section:3.1}

\subsection{Vi điều khiển ESP32}
\label{subsection:3.1.1}

ESP32 là vi điều khiển 32-bit do Espressif Systems phát triển, sử dụng lõi Xtensa LX6 dual-core với xung nhịp tối đa 240\,MHz \cite{espressif_esp32}. Chip tích hợp sẵn Wi-Fi 802.11 b/g/n, Bluetooth 4.2 và BLE, cùng 520\,KB SRAM và hỗ trợ flash ngoài lên đến 16\,MB. Ngoài ra, ESP32 cung cấp đầy đủ giao tiếp ngoại vi như UART, SPI, I\textsuperscript{2}C và ADC, đáp ứng yêu cầu giao tiếp với module SIM7600 qua UART.

Đặc điểm phù hợp với thiết kế này là khả năng giao tiếp đa giao thức ngoại vi: ESP32 sử dụng UART2 để kết nối với module SIM7600, nhận dữ liệu GPS và điều khiển kết nối LTE qua tập lệnh AT command thông qua thư viện TinyGSM. Ngoài ra, ESP32 hỗ trợ ngắt ngoài (external interrupt) trên mọi chân GPIO, cho phép xử lý sự kiện nhấn nút SOS tức thì mà không cần polling liên tục.

Các lựa chọn thay thế đã xem xét gồm: Arduino Uno (tiêu thụ điện cao hơn, thiếu kết nối không dây tích hợp và sức mạnh xử lý không đủ cho TLS), Raspberry Pi Zero W (sử dụng hệ điều hành Linux đầy đủ làm tăng tiêu thụ điện và thời gian khởi động không phù hợp với IoT), và STM32 (phần cứng tốt nhưng hệ sinh thái thư viện cho MQTT và GPS hạn chế hơn). ESP32 được lựa chọn nhờ sự cân bằng tối ưu giữa hiệu năng, tích hợp kết nối không dây, tiêu thụ điện thấp trong chế độ ngủ, và hệ sinh thái Arduino/ESP-IDF phong phú hỗ trợ phát triển nhanh.

\subsection{Module SIM7600}
\label{subsection:3.1.2}

SIM7600 là module đa băng tần LTE Cat-4 do SIMCom Wireless Solutions sản xuất, tích hợp đồng thời bộ thu GNSS (GPS/GLONASS/BeiDou) và modem di động 4G LTE trong một linh kiện duy nhất \cite{simcom_sim7600}. Module giao tiếp với vi điều khiển chủ thông qua giao diện UART sử dụng tập lệnh AT command tiêu chuẩn, cho phép ESP32 điều khiển cả hai chức năng định vị và kết nối dữ liệu qua một kênh giao tiếp duy nhất.

Về hiệu năng định vị, SIM7600 đạt độ chính xác GPS khoảng 2.5\,m CEP (Circular Error Probable) trong điều kiện lý tưởng ngoài trời. Thời gian khởi động lạnh (cold start) khoảng 30 giây, trong khi khởi động nóng (hot start) chỉ khoảng 1 giây khi dữ liệu vệ tinh đã được cache. Firmware giao tiếp với SIM7600 thông qua thư viện TinyGSM — thư viện này trừu tượng hóa toàn bộ tập lệnh AT command, cung cấp API cấp cao như \texttt{modem.getGPS()} để lấy tọa độ và \texttt{modem.gprsConnect()} để kết nối mạng di động, giúp firmware ngắn gọn và không phụ thuộc vào cú pháp AT command cụ thể của từng module.

Các lựa chọn thay thế gồm: ghép module GPS u-blox NEO-6M với modem SIM800L rời (hai module riêng biệt tăng độ phức tạp phần cứng, số lượng linh kiện và diện tích bo mạch), module SIM808 (hỗ trợ GPS và GPRS nhưng không có LTE, băng thông thấp hơn), và GPS tracker module tích hợp sẵn như Quectel EC21 (tương tự nhưng ít tài liệu hỗ trợ hơn trong cộng đồng Việt Nam). SIM7600 được lựa chọn vì tích hợp GPS và kết nối 4G LTE trong một module duy nhất, giảm thiểu số lượng linh kiện, đơn giản hóa firmware, và hỗ trợ tốc độ kết nối cao phù hợp với hạ tầng 4G đang phổ biến tại Việt Nam.

\section{Giao thức truyền thông}
\label{section:3.2}

\subsection{Giao thức MQTT}
\label{subsection:3.2.1}

MQTT (Message Queuing Telemetry Transport) là giao thức nhắn tin nhẹ theo mô hình publish/subscribe, được chuẩn hóa bởi tổ chức OASIS \cite{mqtt_standard}. Giao thức hoạt động ở tầng ứng dụng trên nền TCP/IP với overhead tối thiểu — phần header cố định chỉ 2 byte — phù hợp với kết nối GPRS có băng thông và độ ổn định thấp.

Kiến trúc MQTT dựa trên ba thành phần: \textit{publisher} (thiết bị IoT gửi dữ liệu), \textit{broker} (máy chủ trung gian nhận và phân phối bản tin), và \textit{subscriber} (backend đăng ký nhận bản tin). Thiết bị publish dữ liệu GPS và tín hiệu SOS lên các topic tương ứng; backend đăng ký nhận các topic này bằng cú pháp wildcard và xử lý ngay khi nhận được bản tin. Mô hình này tách biệt hoàn toàn thiết bị và backend, cho phép nhiều subscriber cùng nhận dữ liệu từ một thiết bị mà không cần kết nối trực tiếp.

MQTT hỗ trợ ba mức đảm bảo chất lượng dịch vụ (QoS): QoS 0 (at most once – không đảm bảo), QoS 1 (at least once – đảm bảo giao ít nhất một lần), và QoS 2 (exactly once – đảm bảo giao đúng một lần). Đề tài sử dụng QoS 1 cho các bản tin vị trí: đảm bảo dữ liệu không bị mất trong điều kiện mạng không ổn định với chi phí xử lý thấp hơn QoS 2.

Các giao thức thay thế đã xem xét gồm: HTTP polling (thiết bị phải chủ động gọi API định kỳ, tốn tài nguyên và tăng độ trễ), CoAP (phù hợp UDP nhưng hệ sinh thái thư viện Python hạn chế hơn và ít phổ biến trong IoT di động), và AMQP (mạnh nhưng phức tạp quá mức cần thiết cho ứng dụng này). MQTT được lựa chọn vì là chuẩn công nghiệp cho IoT, có thư viện hỗ trợ tốt trên cả ESP32 (thư viện PubSubClient) lẫn Python backend (thư viện Paho-MQTT \cite{paho_mqtt}), và đặc biệt phù hợp với đặc điểm kết nối LTE của SIM7600.

\subsection{Broker MQTT – EMQX Public Broker}
\label{subsection:3.2.2}

EMQX là MQTT broker hiệu năng cao, mã nguồn mở, được phát triển bởi EMQ Technologies \cite{emqx_broker}. Ngoài phiên bản tự triển khai, EMQ cung cấp dịch vụ broker công khai miễn phí tại địa chỉ \texttt{broker.emqx.io} (cổng 1883 cho kết nối thông thường, cổng 8883 cho TLS), cho phép bất kỳ thiết bị hoặc ứng dụng nào kết nối mà không cần đăng ký tài khoản hay cấu hình xác thực. Broker này hỗ trợ đầy đủ MQTT 3.1, 3.1.1 và 5.0, với khả năng chịu tải lên đến hàng triệu kết nối đồng thời trên hạ tầng production.

Trong hệ thống này, cả thiết bị IoT (ESP32) lẫn backend (FastAPI) đều kết nối đến \texttt{broker.emqx.io} qua cổng 1883. Thiết bị publish dữ liệu GPS và tín hiệu SOS lên các topic riêng biệt; backend đăng ký nhận toàn bộ bằng cú pháp wildcard. Các topic được tổ chức theo cấu trúc phân cấp với namespace dự án riêng, giúp tránh xung đột với các client khác trên cùng broker công khai.

Các lựa chọn thay thế gồm: HiveMQ Cloud (broker quản lý trên đám mây, yêu cầu đăng ký tài khoản và giới hạn kết nối ở gói miễn phí), Eclipse Mosquitto tự triển khai (linh hoạt hơn nhưng yêu cầu quản lý hạ tầng thêm), và AWS IoT Core (chi phí vận hành cao, phức tạp trong cấu hình). EMQX public broker được lựa chọn vì không cần cấu hình thêm, phù hợp cho giai đoạn phát triển và kiểm thử nguyên mẫu, đồng thời dễ dàng thay thế bằng instance EMQX tự triển khai khi đưa vào sản xuất.

\section{Backend}
\label{section:3.3}

\subsection{Framework FastAPI}
\label{subsection:3.3.1}

FastAPI là framework web Python hiện đại, xây dựng trên nền Starlette (ASGI server) và Pydantic (validation dữ liệu), cung cấp hiệu năng cao nhờ hỗ trợ lập trình bất đồng bộ (async/await) \cite{fastapi_docs}. FastAPI tự động sinh tài liệu API theo chuẩn OpenAPI (Swagger UI) từ type hints của Python, giúp kiểm thử và gỡ lỗi trong quá trình phát triển.

Khả năng async của FastAPI đặc biệt phù hợp với yêu cầu của hệ thống này, nơi backend cần xử lý đồng thời nhiều tác vụ I/O-bound: lắng nghe bản tin MQTT từ thiết bị (qua Paho-MQTT callback), quản lý nhiều kết nối WebSocket từ ứng dụng di động, phục vụ yêu cầu REST API, và ghi dữ liệu vào MongoDB. Với mô hình event loop duy nhất, FastAPI xử lý tất cả tác vụ này hiệu quả mà không cần đa luồng phức tạp.

Các lựa chọn thay thế gồm: Django REST Framework (framework đồng bộ, nặng hơn và không tối ưu cho WebSocket và MQTT callback đồng thời), Flask (nhẹ nhưng thiếu hỗ trợ async mạnh, cần thêm thư viện bên ngoài), và Node.js/Express (hệ sinh thái tốt cho realtime nhưng yêu cầu viết lại logic bằng JavaScript thay vì tận dụng hệ sinh thái Python khoa học dữ liệu). FastAPI được lựa chọn nhờ kết hợp tối ưu giữa hiệu năng async, tính đơn giản trong phát triển, và hệ sinh thái Python phong phú.

\subsection{Cơ sở dữ liệu MongoDB}
\label{subsection:3.3.2}

MongoDB là hệ quản trị cơ sở dữ liệu NoSQL hướng tài liệu, lưu trữ dữ liệu dưới định dạng BSON (Binary JSON) \cite{mongodb_docs}. Đối với bài toán lưu trữ dữ liệu vị trí GPS dạng chuỗi thời gian, MongoDB có các ưu điểm: (i) schema linh hoạt phù hợp với payload JSON từ thiết bị IoT, không cần định nghĩa cấu trúc cứng nhắc trước; (ii) Time Series Collection được tối ưu hóa đặc biệt cho dữ liệu theo thời gian, giảm dung lượng lưu trữ và tăng tốc truy vấn theo khoảng thời gian; và (iii) chỉ mục địa lý 2dsphere hỗ trợ truy vấn không gian cho lịch sử hành trình.

Cơ sở dữ liệu gồm sáu collection: \texttt{users} (tài khoản phụ huynh), \texttt{device\_permissions} (liên kết người dùng và thiết bị kèm vai trò), \texttt{device\_configs} (cấu hình thông báo của từng thiết bị), \texttt{locations} (tọa độ GPS theo chuỗi thời gian), \texttt{geofence\_configs} (cấu hình vùng an toàn), và \texttt{alerts} (lịch sử cảnh báo geofence và SOS). Dịch vụ MongoDB Atlas cung cấp gói miễn phí 512\,MB đủ cho giai đoạn thử nghiệm, kèm theo sao lưu tự động và giao diện web quản lý.

Các lựa chọn thay thế gồm: PostgreSQL với PostGIS (mạnh cho dữ liệu không gian nhưng cần cấu hình phức tạp hơn và schema cứng), InfluxDB (chuyên cho time-series nhưng hệ sinh thái Python và driver async hạn chế hơn), và Firebase Realtime Database (dễ tích hợp mobile nhưng chi phí cao khi tải lớn và cấu trúc dữ liệu kém linh hoạt cho truy vấn phức tạp). MongoDB được lựa chọn nhờ tính linh hoạt cao, tích hợp tốt với Python qua thư viện Motor (async driver chính thức), và dịch vụ Atlas miễn phí phù hợp quy mô nguyên mẫu.

\subsection{Xác thực người dùng}
\label{subsection:3.3.3}

Hệ thống xác thực người dùng theo cơ chế JSON Web Token (JWT) kết hợp với bcrypt. Khi đăng ký, mật khẩu được băm bằng thuật toán bcrypt (với salt tự động) trước khi lưu vào cơ sở dữ liệu, đảm bảo không lưu trữ mật khẩu dưới dạng plaintext. Khi đăng nhập, backend xác minh mật khẩu bằng cách so sánh giá trị băm; nếu hợp lệ, server ký và phát hành JWT theo chuẩn HS256 chứa thông tin định danh người dùng và thời hạn hiệu lực. Ứng dụng di động đính kèm token này vào header \texttt{Authorization: Bearer \{token\}} của mọi yêu cầu API tiếp theo; backend xác minh chữ ký và thời hạn token ở mỗi request mà không cần tra cứu session.

Cơ chế JWT stateless này phù hợp với kiến trúc REST và đảm bảo API không thể truy cập bằng token giả mạo hoặc token hết hạn.

\section{Ứng dụng di động}
\label{section:3.4}

\subsection{Framework Flutter}
\label{subsection:3.4.1}

Flutter là framework phát triển ứng dụng đa nền tảng do Google phát triển, sử dụng ngôn ngữ Dart và biên dịch sang mã máy native (AOT compilation) \cite{flutter_docs}. Không giống React Native sử dụng JavaScript bridge để gọi API native, Flutter tự vẽ toàn bộ giao diện thông qua engine đồ họa Skia/Impeller, đảm bảo hiệu năng nhất quán và tốc độ khung hình 60\,fps trên cả hai nền tảng mà không phụ thuộc vào widget của hệ điều hành.

Ứng dụng sử dụng thư viện \textbf{flutter\_map} \cite{flutter_map_pkg} để hiển thị bản đồ tương tác và thư viện \textbf{web\_socket\_channel} để duy trì kết nối WebSocket nhận cập nhật vị trí thời gian thực từ backend. State management được thực hiện bằng Provider pattern, phù hợp với quy mô ứng dụng vừa và nhỏ.

Các lựa chọn thay thế gồm: React Native (hiệu năng thấp hơn với JavaScript bridge và vẫn cần tùy chỉnh native cho một số tính năng bản đồ), phát triển native riêng cho Android (Kotlin) và iOS (Swift) song song (nhân đôi khối lượng công việc phát triển và bảo trì), và Kotlin Multiplatform Mobile (còn đang hoàn thiện, cộng đồng nhỏ hơn). Flutter được lựa chọn nhờ một codebase duy nhất cho cả hai nền tảng, hiệu năng gần native, và hệ sinh thái plugin phong phú đặc biệt cho bản đồ và realtime communication.

\subsection{Bản đồ OpenStreetMap}
\label{subsection:3.4.2}

OpenStreetMap (OSM) là dự án bản đồ địa lý mã nguồn mở do cộng đồng toàn cầu đóng góp, cung cấp dữ liệu bản đồ miễn phí qua định dạng tile ảnh PNG \cite{osm_wiki}. Thư viện flutter\_map tích hợp trực tiếp với OSM tile server, cho phép hiển thị bản đồ tương tác mà không yêu cầu API key hay billing account. Phụ huynh có thể xem vị trí trẻ trên nền bản đồ đường phố với đầy đủ thông tin tên đường, địa danh và công trình.

Lựa chọn thay thế chính là Google Maps SDK: cung cấp bản đồ chất lượng cao và nhiều tính năng nhưng yêu cầu billing account, có giới hạn request miễn phí và chi phí tăng theo lưu lượng sử dụng. Với quy mô nguyên mẫu không xác định được lưu lượng truy cập, OSM qua flutter\_map được lựa chọn để loại bỏ hoàn toàn chi phí và sự phụ thuộc vào dịch vụ của bên thứ ba.

\section{Thuật toán Haversine}
\label{section:3.5}

Thuật toán Haversine tính khoảng cách đường chim bay giữa hai điểm trên bề mặt Trái Đất từ tọa độ kinh độ và vĩ độ của chúng \cite{haversine_formula}. Công thức được biểu diễn như sau:

\begin{equation}
a = \sin^2\!\left(\frac{\Delta\varphi}{2}\right) + \cos\varphi_1 \cdot \cos\varphi_2 \cdot \sin^2\!\left(\frac{\Delta\lambda}{2}\right)
\label{eq:haversine_a}
\end{equation}

\begin{equation}
d = 2R \cdot \arctan2\!\left(\sqrt{a},\, \sqrt{1-a}\right)
\label{eq:haversine_d}
\end{equation}

\noindent trong đó $\varphi_1$, $\varphi_2$ là vĩ độ của điểm đầu và điểm cuối (radian); $\Delta\varphi$ và $\Delta\lambda$ lần lượt là hiệu vĩ độ và kinh độ (radian); $R = 6{,}371$\,km là bán kính trung bình của Trái Đất; và $d$ là khoảng cách cần tính.

Thuật toán Haversine được áp dụng trong backend để kiểm tra vi phạm geofence: mỗi khi nhận được tọa độ mới từ thiết bị, backend tính khoảng cách từ tọa độ đó đến tâm của từng geofence đang hoạt động theo công thức trên. Nếu $d > r$ (bán kính geofence), hệ thống kích hoạt cảnh báo. Đối với chế độ đường đi, thuật toán tính khoảng cách nhỏ nhất từ điểm vị trí đến từng đoạn thẳng của hành trình. Độ phức tạp tổng thể là $O(n \cdot m)$ với $n$ là số geofence và $m$ là số đoạn thẳng trong hành trình — đủ hiệu quả cho quy mô nguyên mẫu.

Các lựa chọn thay thế gồm: Google Geofencing API (phụ thuộc dịch vụ bên ngoài, có chi phí, và không phù hợp với kiến trúc backend tự chủ), PostGIS \texttt{ST\_DWithin} (yêu cầu chuyển sang PostgreSQL), và Vincenty formula (chính xác hơn nhưng phức tạp hơn không cần thiết cho bán kính geofence từ vài chục mét trở lên). Haversine tự triển khai được lựa chọn vì đơn giản, không phụ thuộc dịch vụ ngoài, chính xác đến mức mét và đủ đáp ứng yêu cầu.

Kết chương: Chương 3 đã giới thiệu và phân tích lý do lựa chọn toàn bộ công nghệ trong hệ thống. Sự kết hợp giữa ESP32 và SIM7600 giải quyết yêu cầu về phần cứng nhúng tự chủ với kết nối 4G LTE; MQTT qua EMQX public broker đảm bảo truyền dữ liệu tin cậy trong điều kiện mạng di động; FastAPI và MongoDB cung cấp nền tảng backend async linh hoạt; Flutter và OSM cho phép xây dựng ứng dụng di động đa nền tảng không phụ thuộc dịch vụ có phí; và thuật toán Haversine đáp ứng nhu cầu geofencing hoàn toàn tự chủ. Trên nền tảng công nghệ này, Chương 4 sẽ trình bày chi tiết quá trình thiết kế và triển khai từng thành phần của hệ thống.

\end{document}
